\subsection{I-30. Absorbing endpoint vectors into unit factors; a six-pair width-three sampler}
\label{I-30}

\paragraph{An elementary endpoint factorization.}
Let $u(Y)$ be a polynomial row of length $W$ over a field $L$,
with $u_0=u(0)\ne0$.  Choose $j$ with $(u_0)_j\ne0$ and set
\[
 B(Y)=I_W+\frac{e_j}{(u_0)_j}\bigl(u(Y)-u_0\bigr).
 \tag{I.46}
\]
Then
\[
 u(Y)=u_0B(Y),\qquad B(0)=I_W,\qquad
 \det B(Y)=\frac{u_j(Y)}{(u_0)_j}\ne0.
\]
The first two identities follow by multiplication; the third
is the rank-one determinant identity
$\det(I+ab^T)=1+b^Ta$.
For a polynomial column $v(Z)$ with $v_0=v(0)\ne0$, choosing
$j$ with $(v_0)_j\ne0$ likewise gives
\[
 \begin{gathered}
 C(Z)=I_W+\bigl(v(Z)-v_0\bigr)\frac{e_j^T}{(v_0)_j},
 \qquad v(Z)=C(Z)v_0,\\
 C(0)=I_W,
 \qquad \det C(Z)=\frac{v_j(Z)}{(v_0)_j}\ne0.
 \end{gathered}
 \tag{I.47}
\]
Thus both new endpoint factors are generically invertible and
invertible in every truncated local ring at zero.  No
characteristic restriction is used.

Suppose the interior word has generic interval ranks equal
to its constant interval ranks, in some set of allowed ranks
containing $W$.  Then the enlarged word
$B(Y)A_1\cdots A_m C(Z)$ has the same property.  An interval
lying in the old word is unchanged.  Every other interval is
an old prefix, suffix, or whole word, multiplied on one or
both sides by generically invertible matrices whose constants
are identities; both generic and constant ranks are preserved.
An interval consisting only of a new endpoint factor has
rank $W$.  The same argument holds for interval unit ranks
in a truncated local ring, because multiplication by an
invertible matrix preserves its determinantal rank ideals.

Consequently every coefficient-span support bound proved for
constant endpoints and this rank profile also applies to
\[
 u(Y)A_1\cdots A_m v(Z)
      =u_0B(Y)A_1\cdots A_m C(Z)v_0,
\]
with both endpoint positions included in the original bound.
This costs no additional support positions and is an
existence argument about the representation, not a request
to recover unknown endpoint coefficients from the black box.

\paragraph{Application to all width-three rank tests.}
Use the rank-descent construction of I-21, starting at
$X_i=t$.  Full-rank intervals are preserved as before.
In a selected rank-two test, its second-compound interior
has width three and invertible constant factors.  Its
nonzero polynomial endpoint row and column have nonzero
constants at this center.  Absorb them by (I.46)--(I.47).
The enlarged compound word has all constant factors
invertible, so I-17's row-space argument gives support at
most $3-1=2$ for a nonzero scalar coefficient.  Equivalently
one may use I-28.  Two fresh SV pairs preserve all rank-two
tests simultaneously, since the finite product of their
nonzero substituted minors is nonzero.

The new center is
\[
 \alpha_i=t+\sum_{j=1}^2z_j\lambda_i(y_j).
\]
Every $\alpha_i$ remains transcendental over the original
coefficient field: it has coefficient one on the fresh
indeterminate $t$.  Therefore every nonzero univariate
polynomial endpoint row or column from an original factor
is nonzero at its center.  Constant endpoint vectors are
covered directly.  A selected nonzero contraction cannot
have an identically zero endpoint vector.

In a rank-one test, the interior interval ranks are two or
three and have already been preserved at the current center.
Translate to independent residual variables over the field
of the current parameters.  Absorb both endpoint vectors
using (I.46)--(I.47); every interval of the enlarged word
still has generic and constant rank in $\{2,3\}$.
I-29 with $(w,r)=(3,2)$ now gives support
\[
 h=2(3-2)+3-1=4
\]
for the entire test, including its endpoints.  Four further
fresh SV pairs preserve all the rank-one tests simultaneously.
The preceding rank-two and rank-three minors stay nonzero
because the new amplitudes may be set to zero.
I-21's local-rank induction therefore preserves the rank-one
augmented full word, whose single possibly nonzero entry is
the original scalar $f$.

It follows that the thirteen-parameter map
\[
 \mathcal G_i^{(6)}(t,\mathbf y,\mathbf z)
      =t+\sum_{j=1}^{6}z_j\lambda_i(y_j)
 \tag{I.48}
\]
is a hitting generator for every width-at-most-three ROABP
in one hidden order, in every characteristic.  This is the
same class and evaluation model as I-21; it assumes no
two-rank promise on the whole original word.  The two-rank
lemma is used only inside the rank-one local tests, where
the promise follows from the rank-descent construction.

\paragraph{Exact evaluation bound.}
For total degree at most $D$, the degree in $t$ is at most
$D$, each selector degree is at most $(n-1)D$, and the total
amplitude degree is at most $D$.  Applying the grid and
Newton-simplex argument of I-21 gives the hitting-set bound
\[
 (D+1)\bigl((n-1)D+1\bigr)^6\binom{D+6}{6}.
 \tag{I.49}
\]
The working field must supply at least
$\max\{n,D+1,(n-1)D+1\}$ distinct elements, and evaluations
over that working extension are required if it differs
from the coefficient field.  Earlier ten- and eight-pair
theorems remain valid; (I.48) is the strongest sampler
proved in this track.  Nothing here proves polynomial
dependence on arbitrary width.
\hfill$\square$
